% ==============================================================================
% DANH MỤC TỪ VIẾT TẮT VÀ KÝ HIỆU
% ==============================================================================

\chapter*{DANH MỤC TỪ VIẾT TẮT}
\addcontentsline{toc}{chapter}{Danh mục từ viết tắt}

\begin{table}[h!]
\centering
\renewcommand{\arraystretch}{1.3}
\begin{tabular}{lll}
\toprule
\textbf{Từ viết tắt} & \textbf{Tiếng Anh} & \textbf{Ý nghĩa / Giải thích} \\
\midrule
\textbf{AI} & Artificial Intelligence & Trí tuệ nhân tạo \\
\textbf{AGY} & AlphaGeometry & Kiến trúc chứng minh hình học của DeepMind \\
\textbf{AR} & Algebraic Reasoning & Động cơ suy diễn đại số \\
\textbf{DD} & Deductive Database & Cơ sở dữ liệu suy diễn ký hiệu \\
\textbf{FOL} & First-Order Logic & Logic vị từ bậc nhất \\
\textbf{GraphRAG} & Graph Retrieval-Augmented Generation & Truy xuất tăng cường dựa trên Đồ thị \\
\textbf{IMO} & International Mathematical Olympiad & Kỳ thi Olympic Toán học Quốc tế \\
\textbf{IPS} & Intelligent Problem Solver & Hệ thống giải toán thông minh \\
\textbf{JSON} & JavaScript Object Notation & Định dạng trao đổi dữ liệu tiêu chuẩn \\
\textbf{KB} & Knowledge Base & Cơ sở tri thức \\
\textbf{KG} & Knowledge Graph & Đồ thị tri thức \\
\textbf{KR} & Knowledge Representation & Biểu diễn tri thức \\
\textbf{LLM} & Large Language Model & Mô hình ngôn ngữ lớn \\
\textbf{Neo4j} & Native Graph Database Engine & Cơ sở dữ liệu đồ thị bản địa \\
\textbf{OWL} & Web Ontology Language & Ngôn ngữ biểu diễn Ontology web \\
\textbf{Qdrant} & Vector Similarity Search Engine & Cơ sở dữ liệu tìm kiếm vector tương đồng \\
\textbf{RAG} & Retrieval-Augmented Generation & Tạo sinh tăng cường truy xuất \\
\textbf{REST} & Representational State Transfer & Kiến trúc dịch vụ web RESTful \\
\textbf{SAS} & Side-Angle-Side & Tiên đề bằng nhau Cạnh - Góc - Cạnh \\
\textbf{SSS} & Side-Side-Side & Tiên đề bằng nhau Cạnh - Cạnh - Cạnh \\
\textbf{SSE} & Server-Sent Events & Cơ chế truyền dữ liệu luồng một chiều \\
\textbf{UI} & User Interface & Giao diện người dùng \\
\textbf{URI} & Uniform Resource Identifier & Định danh tài nguyên thống nhất \\
\textbf{WM} & Working Memory & Bộ nhớ làm việc (Tập sự kiện hiện hành) \\
\bottomrule
\end{tabular}
\end{table}

\vspace{0.8cm}

\section*{QUY ƯỚC KÝ HIỆU HÌNH HỌC TRONG TIỂU LUẬN}
\begin{itemize}
    \item $A, B, C, D, \dots$: Các điểm trong mặt phẳng Euclidean $\mathbb{R}^2$.
    \item $AB$ hoặc $\text{Segment}(A, B)$: Đoạn thẳng nối hai điểm $A$ và $B$.
    \item $\text{Length}(AB)$ hoặc $|AB|$: Độ dài hình học của đoạn thẳng $AB$.
    \item $\angle ABC$ hoặc $\text{Angle}(A, B, C)$: Góc đỉnh $B$ tạo bởi hai tia $BA$ và $BC$.
    \item $\triangle ABC$ hoặc $\text{Triangle}(A, B, C)$: Tam giác tạo bởi ba đỉnh $A, B, C$.
    \item $\mathcal{C}(O, R)$ hoặc $\text{Circle}(O)$: Đường tròn tâm $O$ bán kính $R$.
    \item $AB \parallel CD$ hoặc $\text{Parallel}(AB, CD)$: Hai đoạn thẳng/đường thẳng song song.
    \item $AB \perp CD$ hoặc $\text{Perpendicular}(AB, CD)$: Hai đoạn thẳng/đường thẳng vuông góc.
    \item $\triangle ABC \cong \triangle DEF$ hoặc $\text{CongruentTriangles}(ABC, DEF)$: Hai tam giác bằng nhau.
    \item $\triangle ABC \sim \triangle DEF$ hoặc $\text{SimilarTriangles}(ABC, DEF)$: Hai tam giác đồng dạng.
    \item $\mathcal{WM}$: Tập hợp các sự kiện trong bộ nhớ làm việc hiện hành (Working Memory).
    \item $\mathcal{R}$: Tập hợp các luật suy diễn hình học trong cơ sở tri thức (Knowledge Base).
\end{itemize}
